\documentclass{article}
\usepackage[margin=2cm]{geometry}
\usepackage{amsmath}
\usepackage{graphicx}
\usepackage{booktabs}
\usepackage[utf8]{inputenc}
\usepackage{ragged2e}
\setlength{\emergencystretch}{3em}
\begin{document}
\sloppy

e,  e, for all i  j. The free energy in the specific case of a vector multiplet in the
adjoint, a single antisymmetric hypermultiplet, and Ng fundamental hypermultiplets then is.


\begin{equation}
\begin{array} { c } { F ( \lambda _ { i } ) = \sum _ { i \neq j } \left [ F _ { V } ( \lambda _ { i } - \lambda _ { j } ) + F _ { V } ( \lambda _ { i } + \lambda _ { j } ) + F _ { H } ( \lambda _ { i } - \lambda _ { j } ) + F _ { H } ( \lambda _ { i } + \lambda _ { j } ) \right ] } \\ { \qquad + \sum _ { i } \left [ F _ { V } ( 2 \lambda _ { i } ) + F _ { V } ( - 2 \lambda _ { i } ) + N _ { f } F _ {}}\end{array}\right.
\end{equation}


The next step is to look for extrema of this function in the specific case of A,  0 for all.
i. Extrema in the case of non-positive A, can be obtained from these through action of the
Weyl group. To calculate the extrema, one first assumes that as N -> oo, the vevs scale as
X; = Nx, for  > 0 and x of order O(N). One then introduces a density function


\begin{equation}
\rho ( x ) = \frac { 1 } { N } \sum _ { i = 1 } ^ { N } \delta ( x - x _ { i } ) \qquad ( 2 )
\end{equation}


which in the continuum limit should approach an L' function normalized as.


\begin{equation}
\int d x ~ \rho ( x ) = 1 \qquad ( 3 )
\end{equation}


In terms of this density function, one finds that.


\begin{equation}
F \approx - \frac { 9 \pi } { 8 } N ^ { 2 + \alpha } \int d x d y ~ \rho ( x ) \rho ( y ) \left ( | x - y | + | x + y | \right ) + \frac { \pi ( 8 - N _ { f } ) } { 3 } N ^ { 1 + 3 \alpha } \int d x ~ \rho ( x ) ~ | x | ^ { 3 } \quad ( 4 )
\end{equation}


where the large argument expansions have been used, and terms subleading in N have been
dropped. This only has non-trivial saddle points when both terms scale the same with N,
which demands that  = 1/2 and gives the famous result that F x N5/2. Extremizing the
free energy over normalized density functions then gives.


\begin{equation}
F \approx - \frac { 9 \sqrt { 2 } \pi N ^ { 5 / 2 } } { 5 \sqrt { 8 - N _ { f } } } \qquad ( 5 )
\end{equation}


This value of the free energy is to be identified with the renormalized on-shell action of the
supersymmetric AdSe solution. This identification yields the following relation between the
six-dimensional Newton's constant Gs and the parameters N and Ng of the dual SCFT,


\begin{equation}
G _ { 6 } = \frac { 5 \pi \sqrt { 8 - N _ { f } } } { 2 7 \sqrt { 2 } } ~ N ^ { - 5 / 2 } ~
\end{equation}

\end{document}
